% =============================================================================
%  Sección 4: Propuesta de Regeneración
% =============================================================================
\section{Fase 4: Estrategia de Diseño (La Propuesta)}
\label{sec:estrategia-diseño}

La propuesta de intervención del Plan Maestro \textit{Memoria y Función} trasciende 
el simple \termino{embellecimiento} urbano; se constituye como una estrategia integral que 
busca la \textbf{Puesta en Valor} del patrimonio de San Pedro Atocpan mediante 
la solución de sus conflictos de movilidad, seguridad y vulnerabilidad hídrica 
\cite{icomos1964venecia}.

\subsection{Eje de Movilidad Sustentable e Interconexión Regional}
\label{ssec:movilidad-interconexion}

Para garantizar la habitabilidad del centro histórico y la eficiencia de la actividad 
comercial del mole, se propone un modelo de movilidad que prioriza al usuario más 
vulnerable \cite{milpaalta2017pddu}.

\begin{itemize}
    \item \textbf{Centro de Transferencia Multimodal (CETRAM):} Se propone la 
            construcción de un nodo de baja escala colindante con la Carretera Federal 113. 
            Este espacio funcionará como una \termino{rótula de intercambio} donde 
            se interceptarán los autobuses turísticos y de carga, impidiendo su ingreso al 
            Casco Histórico.
    \item \textbf{Dársenas Flexibles y Logística:} El diseño contempla áreas de
            regulación que, durante los fines de semana y la Feria Nacional del Mole,
            operen como estacionamiento disuasorio, permitiendo la redistribución local
            mediante bicitaxis regulados y senderos peatonales \cite{ramirez2017gestion}.
\end{itemize}

La distribución espacial de este sistema multimodal —incluyendo el emplazamiento del
nodo CETRAM, las dársenas flexibles y los flujos de distribución de carga— se
esquematiza en la Figura~\ref{fig:mapa_movilidad_propuesta} del
Anexo~\ref{ssec:planos-fase4}.

\subsection{Rehabilitación Físico-Espacial y Paisaje Urbano Histórico}
\label{ssec:rehabilitacion-paisaje}

La intervención física busca tejer las estructuras urbanas fragmentadas mediante la 
restauración monumental y la recomposición morfológica del tejido 
tradicional \cite{balandrano2016conservacion}.

\begin{enumerate}
    \item \textbf{Tratamiento de Espacios Públicos:} Se plantea la renovación 
            de plazas y calles aplicando un modelo de \textit{sección viaria única} 
            (plataforma compartida) que elimine barreras arquitectónicas y facilite 
            el tránsito peatonal \cite{balandrano2016conservacion}.
    \item \textbf{Recuperación de la Imagen Urbana:} Siguiendo el ejemplo de proyectos 
    exitosos de rehabilitación integral, se implementará un programa de mantenimiento 
    de fachadas corresponsable entre la comunidad y el gobierno, utilizando materiales 
    y colores no invasivos que respeten el paisaje 
    histórico \cite{balandrano2016conservacion, ramirez2017gestion}.
\end{enumerate}

\subsection{Diseño con Perspectiva de Equidad y Seguridad Ciudadana}
\label{ssec:diseño-seguridad}

El espacio público debe planearse para ser seguro y hospitalario para todos los
sectores de la comunidad, evitando que el centro se convierta en una
\termino{escenografía vacía} \cite{rodriguez2009tesis}.

\begin{itemize}
    \item \textbf{Eliminación de Puntos Ciegos:} Se propone la reordenación del 
            mobiliario urbano y la vegetación para mejorar la visibilidad y 
            vigilancia natural, integrando "parques de bolsillo" en zonas 
            actualmente desatendidas.
    \item \textbf{Iluminación Dinámica:} Sistema de alumbrado público que
            no solo garantice la seguridad, sino que enaltezca la calidad
            arquitectónica de la Parroquia de San Pedro y la Capilla de
            San Martín Caballero, facilitando la permanencia social en
            horarios nocturnos \cite{balandrano2016conservacion}.
\end{itemize}

El área de aplicación de estas estrategias corresponde a la Zona Habitacional Intensiva,
cuya delimitación cartográfica —incluyendo los criterios de recuperación de imagen
urbana e iluminación dinámica— se presenta en la
Figura~\ref{fig:mapa_zona_habitacional_intensiva} del Anexo~\ref{ssec:planos-fase4}.

\subsection{Resiliencia Hídrica: Infraestructura Verde y Renaturalización}
\label{ssec:resilencia-infraestructura}

Dada la etimología de Atocpan y el precedente del desastre de 1935, la 
propuesta integra la variable ambiental como eje de competitividad y 
seguridad \cite{balandrano2016conservacion}.

\begin{itemize}
    \item \textbf{Renaturalización Urbana:} Conexión del paisaje natural del 
            Chichinautzin con el paisaje construido mediante la creación de 
            corredores verdes y superficies permeables 
            (pavimentos pétreos filtrantes) que reduzcan el 
            escurrimiento superficial \cite{balandrano2016conservacion, rodriguez2009tesis}.
    \item \textbf{Sistemas de Captación Pluvial:} Instalación de infraestructura
            para el manejo y aprovechamiento de tecnologías limpias en
            edificios públicos, compensando hundimientos diferenciales y
            evitando la saturación del sistema de
            drenaje inconcluso \cite{balandrano2016conservacion}.
\end{itemize}

Las áreas de intervención ambiental y las zonas de restricción por vulnerabilidad
hídrica se delimitan en las Figuras~\ref{fig:mapa_zona_habitacional_extensiva}
(Anexo~\ref{ssec:planos-fase4}) y~\ref{fig:mapa_zona_riesgo}
(Anexo~\ref{ssec:planos-zonificacion}).